\section{电场环路定理与电势}
在上一节中我们研究了电场强度在曲面上的积累，即电通量，进而我们通过电场强度在闭曲面上的积分，得出了\textbf{电场的高斯定理}和\textbf{散度性质}。类似的，本节将研究电场强度在曲线上的积累，即电势差，进而通过电场强度在闭曲线上的积分，得出\textbf{电场的环路定理}和\textbf{旋度性质}。

% 在本节中我们将研究电场强度在曲线上的积累，即电势，电场强度反映的是电场“力”的性质，电势则反映电场“能”的性质

\subsection{电势能}
电势是电场强度在空间上的累积，为了研究电势，我们就要先研究电场力在空间上的累积，即电场力的做功的性质，本小节将由此引出电势能的概念。

根据\Refx{定律}{库仑定律}
\begin{equation*}
    \vb*{F}=\frac{1}{4\pi\varepsilon_0}\frac{qq_0}{r^2}\vu*{r}
\end{equation*}
其中$\vb*{F}$代表试探电荷$q_0$受到的由$q$产生的电场的电场力，而$\vb*{r}$代表$q_0$的位矢
\begin{equation*}
    \vb*{F}=\frac{1}{4\pi\varepsilon_0}\qty(\frac{qq_0x}{(x^2+y^2+z^2)^\frac{3}{2}}\vi-\frac{qq_0y}{(x^2+y^2+z^2)^\frac{3}{2}}\vj-\frac{qq_0z}{(x^2+y^2+z^2)^\frac{3}{2}}\vk)
\end{equation*}
关注到偏导数（另外两组偏导数也可以类似求得且也是相等的）
\begin{align*}
    \Fpif{F_x}{y}&=\frac{1}{4\pi\varepsilon_0}\frac{qq_0x\cdot\dfrac{3}{2}(x^2+y^2+z^2)^{\frac{1}{2}}\cdot 2y}{(x^2+y^2+z^2)^{\frac{3}{2}}}=\frac{1}{4\pi\varepsilon_0}\frac{3qq_0xy}{x^2+y^2+z^2}\\[3mm]
    \Fpif{F_y}{x}&=\frac{1}{4\pi\varepsilon_0}\frac{qq_0y\cdot\dfrac{3}{2}(x^2+y^2+z^2)^{\frac{1}{2}}\cdot 2x}{(x^2+y^2+z^2)^{\frac{3}{2}}}=\frac{1}{4\pi\varepsilon_0}\frac{3qq_0yx}{x^2+y^2+z^2}
\end{align*}

这就满足空间曲线积与路径无关的条件
\begin{equation*}
    \Fpif{F_x}{y}=\Fpif{F_y}{x}\qquad\Fpif{F_y}{z}=\Fpif{F_z}{y}\qquad\Fpif{F_z}{x}=\Fpif{F_x}{z}
\end{equation*}
因此电场力做功与路径无关，只与始末位置有关，即\textbf{电场力是一种保守力}。

这就使得对于电场中，由点$A$至点$B$的任意路径上电场力对试探电荷$q_0$所做的功，可以等效的化为两段路径（假定$r_a<r_b$，其中$r_a,r_b$分别的代表点$A,B$至场源电荷$q$的距离）
\begin{enumerate}
    \item 第一段由$A$沿极径到达与$B$半径相同的$B_0$处，这是一个远离场源电荷$q$的过程，因此，若$qq_0>0$则电场力与位移方向相同，若果$qq_0<0$则电场力与位移方向相反。
    \item 第二段由$B_0$沿球面到达$B$处，该过程中电场力总是垂直于球面，故不做功。
\end{enumerate}
因此电场力做功可以转化为关于$r$的定积分
\begin{Equation}[电势能的准备式]
    W=\frac{1}{4\pi\varepsilon_0}\Int[r_a][r_b]{\frac{qq_0}{r^2}\dif{r}}=\frac{1}{4\pi\varepsilon_0}\qty(-\frac{qq_0}{r_b})-\qty(-\frac{qq_0}{r_a})
\end{Equation}
如果我们引入\uwave{电势能}的概念
\begin{BoxDefinition}[电势能]
    E_p=\frac{1}{4\pi\varepsilon_0}\frac{qq_0}{r}
\end{BoxDefinition}
那么式\ref{式:电势能的准备式}即$W=(-E_{pb})-(-E_{pa})=-\delt{E_p}$，即\textbf{电场力所做的功是电势能的负增量}。

\Refx{定义}{电势能}指出，试探电荷在点电荷激发的电场中所具有的电势能，与试探电荷和场源电荷的电荷量乘积成正比，与到试探电荷到场源电荷的距离成反比。

\Refx{定义}{电势能}与\Refx{定律}{库仑定律}的公式在形式上非常接近，但应注意两点区别，电场力是矢量且大小按照距离的平方反比减小，电势能是标量且大小按照距离的反比减小。

以上推导和结论与先前力学部分的\Refx{定义}{引力势能}十分相似，这是因为电场力和引力都是平方反比力，区别仅在于前者是质量乘积后者是电荷量乘积，两者力的相似数学形式决定了两者对应势能的相似数学形式，即电势能和引力势能均以距离反比衰减。

电势能与我们熟知的势能具有相同的性质
\begin{itemize}
    \item 当电势能做正功时，电荷的电势能减少，动能增加。
    \item 当电势能做负功时，电荷的电势能增加，动能减少。
\end{itemize}

通过观察不难发现，对于点电荷的电势能，当$r\rightarrow +\infty$时有$E_{p}\rightarrow 0$，因此点电荷的电势能实际是以无穷远处为势能零点，即电场中$A$点处的电势能可以表示为
\begin{BoxEquation}[电势能的积分表示]
    E_p=\Int[A][+\infty]{\vb*{F}\cdot\dif{\vb*{r}}}
\end{BoxEquation}

\subsection{电势}
电势能$E_p=\dfrac{1}{4\pi\varepsilon_0}\dfrac{qq_0}{r}$的大小与试探电荷$q_0$的电荷量多少有关，因此无法客观的反映电场本身的性质，但关注到比值$\dfrac{E_p}{q_0}=\dfrac{1}{4\pi\varepsilon_0}\dfrac{q}{r}$与试探电荷$q_0$无关，故$\dfrac{E_p}{q_0}$可以反映电场性质。\vspace{8pt}

定义单位正试探电荷在电场中具有的电势能，为该点处的\uwave{电势}，使用$V$表示，即
\begin{BoxDefinition}[电势]
    V=\frac{E_p}{q_0}
\end{BoxDefinition}

\Refx{定义}{电势}中代入\Refx{定义}{电势能}，就可以得到点电荷的电势计算公式
\begin{BoxEquation}[点电荷的电势]
    V=\frac{1}{4\pi\varepsilon_0}\frac{q}{r}
\end{BoxEquation}
\Refx{公式}{点电荷的电势}指出了点电荷的电势分布规律
\begin{itemize}
    \item 正点电荷周围的电势均为正，且靠近正点电荷处的电势具有较大的正值。
    \item 负点电荷周围的电势均为负，且靠近负点电荷处的电势具有较大的负值。
\end{itemize}

\Refx{定义}{电场强度}和\Refx{定义}{电势}中，我们分别用单位正试探电荷在电场中受到的电场力和具有的电势能，定义了电场强度和电势的概念，两者从不同的方面反映了电场具有的性质——电场强度反映了电场“力的性质”，电势反映了电场“能的性质”。

引入电势的概念后，电场中的点电荷所受的电场力就可以表示为
\begin{BoxEquation}[电势能的电势表示]
    E_p=Vq_0
\end{BoxEquation}
\Refx{公式}{电势能的电势表示}说明
\begin{itemize}
    \item 正点电荷在电场中具有的电势能，与该点处的电势符号相同。
    \item 负点电荷在电场中具有的电势能，与该店处的电势符号相反。
\end{itemize}

电势也可以以积分形式表示，\Refx{公式}{电势能的积分表示}中代入$\vb*{F}=\vb*{E}q_0$和$\vb*{E}_p=Vq_0$
\begin{BoxEquation}[电势的积分表示]
    V=\Int[A][+\infty]{\vb*{E}\cdot\dif\vb*{r}}
\end{BoxEquation}
\Refx{公式}{电势的积分表示}指出，某点处的电势是电场强度由该点至参考点的第二类曲线积分，具体路径无关紧要，电场力的积分是路径无关，电场强度的积分也是路径无关的。

对于点电荷的情况，电势的参考点，即零电势点也常常选在无穷远处，这对于大部分的点电荷系和电荷连续分布体来说都是没有问题的，唯一需要特别注意的是，当我们在尝试研究类似于“无限大带电平板”这类自身线度就为无穷大的对象时，我们就不宜再选取无穷远为零电势点了，因为此时的无穷远并非各向同性的，这时通常会取带电体本身为零电势点。

\Refx{公式}{电势的积分表示}指出，电场强度是电势的负梯度
\begin{BoxEquation}[电场强度与电势的梯度关系]
    \vb*{E}=-\nabla V
\end{BoxEquation}
\Refx{公式}{电场强度与电势的梯度关系}就表明了电场的一个重要性质
\begin{center}
    电场是保守场
\end{center}
电势在国际单位制中的单位是\si{V=J\cdot C^{-1}}，称为\uwave{伏特}，量纲是\si{[L^2MT^{-3}I^{-1}]}。

\subsection{电势叠加原理}
当空间中存在多个场源电荷$q_1,q_2,\cdots,q_n$时，在场点$P$处的电场强度为$\vb*{E}_1,\vb*{E}_2,\cdots,\vb*{E}_n$，而电场强度是可以合成的，因此场点$P$处的电场强度$\vb*{E}=\vb*{E}_1+\vb*{E}_2+\cdots+\vb*{E}_n$。

根据\Refx{定义}{电势}，点$P$处的场强$\vb*{E}$
\begin{equation*}
    \vb*{V}=\Int[A][+\infty]{\vb*{E}\cdot\dif{\vb*{r}}}=\Int[A][+\infty]{(\vb*{E}_1+\vb*{E}_2+\cdots+\vb*{E}_n)}=V_1+V_2+\cdots+V_n
\end{equation*}
其中$V_1,V_2,\cdots,V_n$表示各场源电荷$q_1,q_2,\cdots,q_n$独立存在时产生的电势，因此
\begin{BoxPrinciple}[电势的叠加原理]
    点电荷系在电场中某点处所产生的电势，
    
    等于各点电荷单独存在时，在该点处所产生的电势度的代数和。
    \begin{BoxEq}
        V=\SumIN{V_i}
    \end{BoxEq}
\end{BoxPrinciple}
电荷连续分布的任意带电体产生的电势，可以将带电体携带的电荷看作许多电荷元$\dif{q}$组成，每个电荷元$\dif{q}$均可以视为点电荷，其在距离$r$处产生的电势$\dif{V}$为
\begin{BoxEquation}[电荷元的电势]
    \dif{V}=\frac{1}{4\pi\varepsilon_0}\frac{\dif{q}}{r}
\end{BoxEquation}
电荷连续分布的带电体产生的电势，需要通过对$\dif{V}$积分完成，积分形式取决于电荷分布，积分的具体形式可以参照\Refx{公式}{电荷元的电场强度}处的公式，这里不再赘述。

\subsubsection{电偶极子的电势}
现在我们来研究，电偶极子远场点处的电势。

设远场点$P$相对于电偶极子中心的位矢为$\vb*{r}$，至正负电荷的距离为$r_{+},r_{-}$。
\begin{Figure}{电偶极子的电势}
    \inputfigure{Chapter08C_Figure01}\vspace{-10pt}
\end{Figure}
根据\Refx{公式}{点电荷的电势}，$\pm q$单独存在时在点$P$产生的电势为
\begin{Equation}[电偶极子的电势]
    V_{+}=\ek\frac{q}{r_{+}}\qquad
    V_{-}=\ek\frac{q}{r_{-}}
\end{Equation}
过$\pm q$所在的点$A,B$向$PO$做垂线交其于垂足$H_{+},H_{-}$，容易看出$OH_{+}=OH_{-}=\dfrac{l}{2}\cos\theta$，其中$\theta$为矢量$\vb*{l},\vb*{r}$夹角，由于这里$P$为远场点，这里近似认为$PH_{+}\approx PA$和$PH_{-}\approx PB$。

从而$r_{+}$和$r_{-}$均可以统一的用$r$来表示
\begin{equation*}
    r_{+}=r-\frac{l}{2}\cos\theta\qquad
    r_{-}=r+\frac{l}{2}\cos\theta
\end{equation*}
将上式代入\ref{式:电偶极子的电势}，得到
\begin{align*}
    V=V_{+}+V_{-}=\ek\qty(\frac{q}{r-\frac{l}{2}\cos\theta}-\frac{q}{r+\frac{l}{2}\cos\theta})=\ek\qty(\frac{ql\cos\theta}{r^2-\qty(\frac{l}{2}\cos\theta)^2})
\end{align*}
由于$r\gg l$，因此上式分母中的$\qty(\frac{l}{2}\cos\theta)$可以略去
\begin{equation*}
    V=\ek\frac{ql\cos\theta}{r^2}=\ek\frac{p\cos\theta}{r^2}
\end{equation*}
将上式通过矢量的形式表示，即
\begin{BoxEquation}[电偶极子的电势]
    V=\ek\frac{\vb*{p}\cdot\vu*{r}}{r^2}
\end{BoxEquation}

\subsection{电势差}
电势能的负增量是电场力所做的功，电势与电势能相对应，电势的负增量也应有一个对应的物理量，我们将电势的负增量定义为\uwave{电势差}，也称为\uwave{电压}，用符号$U$表示
\begin{BoxDefinition}[电势差]
    U_{ab}=V_a-V_b=-\delt{V}
\end{BoxDefinition}
这里应特别注意，电势的增量$\delt{V}=V_b-V_a$是末减初，电势差$U_{ab}=V_b-V_a$是初减末，
两者是相反的，\textbf{电势差和电势的增量是不同的概念}，\textbf{电势差被定义为电势的负增量}，这可能听上去十分别扭，但是如果将电场力的功$W_{ab}=(-E_{pb})-(-E_{pa})=-\delt{E}_p$展开
\begin{BoxEquation}[电场力的功]
    W_{ab}=E_{pa}-E_{pb}=-\delt{E_{p}}
\end{BoxEquation}
这样形式就对应了，电场力的功$W_{ab}$是电势能的负增量，电势差$U_{ab}$是电势的负增量，根据\Refx{公式}{电势能的电势表示}可知$E_p=Vq_0$，即$\delt{E_p}=\delt{V}q_0$，故有
\begin{BoxEquation}[电势差与电场力的功的关系]
    U_{ab}=\frac{W_{ab}}{q_0}
\end{BoxEquation}
\Refx{公式}{电势差与电场力的功的关系}指出，在$A,B$两点间的电势差$U_{ab}$，等于将单位正电荷由$A$移动到$B$时电场力所做的功$W_{ab}$，由此可见“电势差与电场力的功”的关系，与先前“电场强度和电场力”以及“电势与电势能”的关系一样，后者是前者的$q_0$倍。

电压和电势差在国际单位制中的单位与电势相同，均为\si{V=J\cdot C^{-1}}。

根据\Refx{定义}{变力的功}，电场力所做的功可以用积分表示为
\begin{BoxEquation}[电场力的功的积分表示]
    W_{ab}=\Int[A][B]{\vb*{F}\cdot\dif{\vb*{r}}}
\end{BoxEquation}
根据\Refx{公式}{电势差与电场力的功的关系}，电势差可以对应的用积分表示为
\begin{BoxEquation}[电势差的积分表示]
    U_{ab}=\Int[A][B]{\vb*{E}\cdot\dif{\vb*{r}}}
\end{BoxEquation}
\Refx{公式}{电势差的积分表示}的一个特例是在匀强电场中
\begin{Figure}{匀强电场中的电势差}
    \inputfigure[scale=0.8]{Chapter08C_Figure02}
\end{Figure}
\Refx{公式}{电势差的积分表示}此时可以表示为
\begin{BoxEquation}[匀强电场的电势差]
    U_{ab}=\vb*{E}\cdot\delt\vb*{r}    
\end{BoxEquation}
其中$\delt{\vb*{r}}$表示由点$A$指向点$B$的位移矢量。

\subsection{等势面}
为了形象描绘电场的电势分布，我们引入等势面的概念，我们将空间中电势相等的点所构成的曲面称为\uwave{等势面}，如果我们等距的取一系列电势值，那么就会得到一簇等势面，例如
\begin{itemize}
    \item 孤立点电荷的等势面是一簇以点电荷为中心的同心球面，且越靠近电荷等势面越密集。
    \item 带电平板间的匀强电场的等势面是一簇空间上等距分布的平面。
\end{itemize}
我们知道，电场线通常以带箭头的直线绘制，等势面则通常以虚线绘制，因为通常我们绘制的电场分布图都是某一截面上的情况，因此绘出的虚线实质是等势面在观察截面上的截线。

\Refx{图}{匀强电场中的电势差}展示了匀强电场的电场线和等势面的分布，点电荷与点电荷系的等势面分布参见\Refx{表}{常见带电系统的电场分布}中的彩色虚线，其颜色表征了电势的高低，越接近暖色（红色）表明电势正值越高，越接近冷色（紫色）表明电势负值越高。

等势面和电场线具有以下关系
\begin{itemize}
    \item 电场线总是与等势面垂直，电场强度越大处的等势面越密集。
    \item 电场线指向电势下降最快的方向，顺电场线方向电势降低，逆电场线方向电势升高。
    \item 电荷在同一等势面上运动时，电场力不做功，电势能不改变。
\end{itemize}
等势面和电场的关系是由\Refx{公式}{电场强度与电势的梯度关系}决定的。\vspace{-3pt}

\subsection{电场环路定理}
电场是保守场，因此电场强度沿任意空间闭合曲线的环路积分均为零，即
\begin{BoxPrinciple}[电场环路定理]
    在电场中，沿着任意一个闭合曲线的电场强度环流均为零。
    \begin{BoxEq}
        \Ilot[\Gamma]{\vb*{E}\cdot\dif{\vb*{r}}}=0
    \end{BoxEq}
\end{BoxPrinciple}
该规律就称为\uwave{电场环路定理}。

根据旋度的定义，电场在某一点处的旋度，是该点处的电场强度的环流密度，即\footnote{其中$P$为电场中某一点，而$S$表示闭曲线$\Gamma$所张成的空间曲面的面积。}
\begin{equation*}
    \nabla\times\vb*{E}=\Lim{\Gamma}{P}{\frac{1}{S}\Ilot[\Gamma]{\vb*{E}\cdot\dif{\vb*{r}}}}=\vb*{0}
\end{equation*}
这就得到了一个重要定律
\begin{BoxPrinciple}[电场环路定理的微分形式]
    在电场中，电场强度的旋度恒为零
    \begin{BoxEq}
        \nabla\times\vb*{E}=\vb*{0}
    \end{BoxEq}
\end{BoxPrinciple}
\Refx{定律}{电场环路定理的微分形式}也可以从另外一个角度得出，先前我们已经知道电场是保守场，即电场是某个标量场的梯度场，而梯度场必然无旋，这就有$\nabla\times\vb*{E}=\vb*{0}$。

\Refx{定律}{电场高斯定理的微分形式}指出
\begin{center}
    电场是无旋场
\end{center}
电场是无旋场，这一点也反映在电场线的特征上——电场线不会形成闭合曲线。

至此，我们可以归纳出以下的电场基本性质
\begin{BoxPrinciple}[电场的基本性质]
    电场是一个有源无旋的保守场。
\end{BoxPrinciple}
